\chapter{Literature Review}\label{chap:background}


\section{Emotion in Learning}\label{sec:emotion_in_learning}

Early work in educational psychology treated emotion in achievement settings mainly through the lens of test anxiety. Foundational studies linked anxiety to interference with learning and performance in evaluative situations \cite{mandler1952anxiety}, and later reviews documented substantial correlational and experimental evidence on how anxiety relates to cognition and academic performance \cite{hembree1988testanxiety,zeidner1998testanxiety}. Even within that limited focus, researchers were already trying to understand achievement-related emotions through the way students explain success and failure \cite{weiner1985attributional}. What remained comparatively open, however, was whether anxiety was the only emotion that mattered for how students actually learn in classrooms and during study.

By the early 2000s, that gap was addressed directly. Qualitative and questionnaire-based work showed that students report a wide range of emotions tied to instruction, studying, and testing, not only fear around exams \cite{pekrun_academic_nodate}. The same line of work introduced the term \emph{academic emotions} for this wider group and classified them using two simple dimensions: whether the emotion is positive or negative, and whether it increases or reduces a student's level of activity. Emotions were understood to influence learning and performance indirectly through motivation, learning strategies, mental resources, and self-regulation, while students' results and feedback could in turn influence how they feel in future learning situations \cite{pekrun_academic_nodate}. In other words, researchers increasingly came to understand that learning is affected not only by one negative emotion, but by a broader set of emotional states that may appear during the same lesson and shape the way students pay attention, stay involved, participate, and make progress.

A parallel line of research developed around the concept of engagement rather than focusing only on specific emotion labels, although it still included a clear emotional dimension. In this view, engagement was not treated as a single construct, but as a combination of three related aspects: behavioural participation in learning activities, emotional reactions to school and instruction, and cognitive effort directed towards understanding difficult material \cite{fredricks_httpwwwjstororg_2004}. Emotional engagement in that sense covers interest, boredom, anxiety, and related feelings towards school and academic work, and it is treated as part of why students are willing to keep trying when tasks are hard \cite{fredricks_httpwwwjstororg_2004}. Fredricks and colleagues connected these dimensions of engagement to long-standing motivational theories which suggest that students are more likely to remain involved in learning when they feel competent, have a sense of control over their learning, and feel socially connected to others \cite{connell1991competence}. They also linked them to expectancy-value perspectives, which suggest that students' decisions to engage are shaped by how interesting they find a task and how important or useful they believe it to be \cite{eccles1983expectations}. Engagement therefore became an additional way of understanding affect in education. It did not replace the study of specific emotions, but worked alongside it by offering a broader view of students' emotional involvement in learning.

From the mid-2000s onward, these educational perspectives were strengthened by more explicit theories explaining how emotions arise in learning situations, as well as by neuroscientific work on why affect matters for learning in the first place. One important example is Pekrun's control-value theory of achievement emotions, which argues that students' emotions depend largely on how much control they believe they have over academic tasks and how much value they attach to those tasks. According to this view, different combinations of control and value can give rise to emotions such as enjoyment, hope, anxiety, boredom, and shame, and these emotions can then influence students' interest, motivation, engagement, and academic achievement \cite{pekrun_control-value_2006}. Independently, Immordino-Yang and Damasio argued that higher-level learning in school is not separate from emotion. They described emotional thought as the connection between feelings and the attention, memory, and reasoning processes that education tries to strengthen, and explained that emotion helps make experiences meaningful, which allows knowledge to be remembered and later used in real decisions \cite{immordino-yang_we_2007}. This argument builds on earlier work which suggested that emotion is an essential part of rational behaviour and decision-making, rather than something that simply distracts from clear thinking \cite{damasio1994descartes}. Together, the appraisal perspective and the neuroscience perspective support the same general idea: emotion is not just an extra layer added to learning, but part of the process itself. In this way, these lines of work helped move educational research beyond a purely cognitive view of learning towards one that also recognises the role of motivation and affect.

Research on dialogue-based tutoring then helped explain more clearly how emotions appear during learning, especially when students interact with a tutor while solving problems, asking questions, and receiving feedback. Studies of long tutoring conversations showed that deep learning is often accompanied by a range of emotions, especially boredom, engagement, confusion, and frustration, which appear more often than the traditional basic emotion categories would suggest \cite{dmello_language_2012}. These four states later became the de facto target set in webcam-based educational facial expression recognition benchmarks such as \ac{DAiSEE}, which the following sections review. Experimental studies showed that a moderate level of confusion can sometimes improve learning outcomes, partly because it increases attention and encourages deeper processing of the material \cite{dmello2014confusionbeneficial}. A later review by Tyng and colleagues connected these findings from classrooms with brain imaging studies that showed how emotion affects thinking and learning \cite{tyng_influences_2017}. Information with emotional meaning usually attracts attention more quickly and is given more focus than neutral information \cite{schupp2007selective,vuilleumier2005brainsbeware}. At the brain level, areas such as the amygdala, medial temporal lobe, and prefrontal cortex help explain how emotional arousal influences the formation and strengthening of memories, which is why emotional events are often remembered more clearly than neutral ones \cite{tyng_influences_2017,sharot2004emotionremembering}. The same review stresses that stress and arousal can help or hurt learning depending on intensity and timing, and that motivational states such as curiosity support exploration and readiness to learn.

More recent review studies still show that emotion is a central part of how students experience learning, especially in difficult subjects and in online settings where students often receive less informal support. A systematic review of emotions in STEM online learning found that emotions strongly influence engagement, motivation, and satisfaction, and that teaching needs a better understanding of students' emotional experiences in order to respond to what they are actually facing \cite{anwar_emotions_2023}. This conclusion builds on a long tradition of educational psychology research that had already treated emotion as an important topic to understand and measure in schools \cite{schutz2007emotioneducation}. Once emotion is treated as part of how students attend, remember, and stay engaged, the practical question becomes how to tell what a learner is feeling while they are actually working. The next section therefore turns from the conceptual case for emotion in learning to the ways researchers have tried to observe or measure it in educational settings.


\section{Emotion Detection in Educational Systems}\label{sec:emotion_detection_in_educational_systems}

Once emotions were established as central to learning, motivation, and achievement, the field faced a more operational question: how can those emotions be detected while learning is actually taking place. Early work in affective computing had already described emotion detection as a measurement problem that could use several different signals rather than only one method \cite{Picard1997AffectiveComputing}, and educational researchers gradually translated that idea into classroom and tutoring contexts \cite{Calvo2010AffectDetection}. What came next was not one simple shift from one method to another, but several overlapping approaches, each with different strengths and limitations in terms of timing, scalability, and practical use in real educational settings.

At first, most educational studies relied on self-report. The achievement-emotions line of work around the \ac{AEQ} gave researchers a structured way to measure recurring states such as enjoyment, pride, anxiety, anger, shame, and boredom in relation to classroom activity and assessment \cite{pekrun_control-value_2006,pekrun_academic_nodate}. The same type of questionnaires is still used in recent classroom AI studies. For example, Omar and colleagues used \ac{AEQ-M} surveys before and after the intervention, along with interviews, to track emotional changes in grade 8 algebra classes using an AI-supported learning environment \cite{omar_academic_2025}. Studies in distance learning also refined self-report methods instead of abandoning them. For example, some studies compared daily questionnaires completed after learning activities with emoticon tools used during the course, where students could mark moments when they felt certain emotions. These methods had clear advantages, including low instrumentation demands, direct access to students' own experiences, and easy deployment across different settings. However, they also had persistent limitations, such as dependence on students' self-awareness, delayed reporting, sensitivity to wording and social context, and difficulty capturing rapid emotional changes in real time \cite{Lavoue2017EmotionalDataCollection}.

Alongside these questionnaire-based approaches, researchers also used direct human observation to understand students' affect and engagement from their visible behaviour in real classrooms. In school psychology, structured coding systems such as \ac{AET-SSBD}, \ac{BOSS}, \ac{DOF}, \ac{COC}, \ac{SECOS}, \ac{ADHD-SOC}, and \ac{SOS} were used to observe students at set time intervals and record whether they were focused, distracted, or showing signs of difficulty \cite{Volpe2005ObservingStudents,Shapiro2004AcademicSkillsProblemsWorkbook}. Related work in learning technologies also used a similar observation-based approach in educational software and classroom interactions \cite{Baker2010BetterFrustratedThanBored,Liu2017DetectStudentsAcademicEmotions}. Compared with questionnaires completed after the event, observation methods gave more detailed information about when emotions and behaviours appeared and linked interpretations more closely to what students were actually doing in the learning setting. However, these methods were demanding because they required observer training, checks for consistency between observers, repeated observations over time, and careful control of problems such as observer drift and reactivity, while still showing only small parts of students' experience during the full course \cite{Volpe2005ObservingStudents}.

As digital tutoring systems collected detailed data about what students did, researchers began to move from manual coding to large-scale behavioural inference. Instead of asking students about their feelings or observing them all the time, they used records such as answer correctness, hint use, timing, and help-seeking patterns to infer affect and disengagement \cite{Baker2007OffTaskBehaviorITS,Cocea2009OffTaskGamingLearning}. Pardos and colleagues are a good example of this shift. They built detectors in \ac{ASSISTments} that did not need sensors and were used to identify states such as boredom, engaged concentration, confusion, frustration, off-task behaviour, and gaming the system. They then followed these affective patterns over a full year and examined how they were related to students' end-of-year mathematics test results \cite{Pardos2014AffectiveStatesStateTests}. Studies in technology-rich classrooms also used coded behaviour data, sequence analysis, and Hidden Markov Models to study how students' emotions and behaviours changed over time \cite{Liu2017DetectStudentsAcademicEmotions}. These approaches made it easier to study many students and to follow them over time, but they still estimated emotion indirectly from behaviour. This is a limitation because the same behaviour can reflect different emotions, the models depend on labelled training data, and their performance may vary across platforms, tasks, and teaching contexts \cite{Pardos2014AffectiveStatesStateTests}. The same behavioural traces later became the default input to many adaptive and reinforcement-learning tutors, even when facial or other perceptual signals were available \cite{zerkouk_comprehensive_2025,meng_self-evolving_2025}.

Dialogue-based tutoring introduced another approach by using language and conversation as signals of emotion. Research on tutoring dialogue showed that emotional states could, in principle, be inferred from what learners and tutors said, even when explicit emotion words were not used \cite{DMello2008AutomaticDetectionLearnersAffect}. In typed \ac{AutoTutor} dialogues, D'Mello and Graesser demonstrated that discourse-sensitive features derived from \ac{LIWC} and \ac{Coh-Metrix} could predict distributions of boredom, confusion, flow/engagement, and frustration over whole sessions \cite{dmello_language_2012,Graesser2004AutoTutorDialogue}. Those results showed that language carries affective information, but only under conditions that are difficult to reproduce in ordinary online study. The models depend on extended tutoring conversation, domain-specific vocabulary, and session-level transcript summaries rather than on the short, ambiguous messages that dominate most digital exercises. A student who types ``I don't get it'', ``ok'', or nothing at all may be confused, bored, or simply unwilling to write; the same string of words can therefore support more than one emotional reading, which makes text an inherently ambiguous channel for real-time adaptation \cite{dmello_language_2012,Lavoue2017EmotionalDataCollection}. Recent chat-based and large-language-model tutors inherit the same problem: they can process what the learner chooses to write, but they cannot see the frustration or disengagement that appears on the face when the learner stops typing \cite{gutierrez_development_2025,meng_self-evolving_2025}. Surveys of sentiment analysis in STEM online learning report the same weakness at scale: text-based methods struggle with sparse emotional language, context dependence, and the gap between what students write and what they actually feel during a task \cite{anwar_emotions_2023,florestiyanto_emotion_2024}.

Facial expression is a stronger and more direct alternative for educational affect monitoring. Comparative work during classroom use of an educational game found that face-only affect detectors were more accurate than detectors built from interaction logs alone, even though facial models could not be applied in every student case \cite{bosch_2015}. That pattern is consistent with a wider lesson from the field: visible affect on the face is closer to the emotional event than inferred labels derived from clicks, timings, or typed responses. Teachers already rely on this signal in classrooms; once webcams became standard in online learning, the same cue could be captured continuously without asking students to describe their feelings in words. Even early affective tutoring systems that experimented with text semantics treated facial expression as a necessary companion rather than an optional extra, because text alone did not provide a stable enough reading of the user's state \cite{lin_employing_2012}. For adaptive systems that must respond while a student is still working, facial behaviour therefore offers a clearer, more continuous, and less ambiguous basis for detecting academic emotions such as confusion, boredom, frustration, and engagement than dialogue or log data alone \cite{lek_academic_2023,vistorte_integrating_2024,petrovica_emotion_2016}. The next section reviews how educational facial expression recognition has developed around that advantage.

\section{Facial Expression Recognition in Education}\label{sec:fer_in_education}

Among the signals reviewed in the previous section, one stood out as visible and easy to capture during everyday learning. A learner's face is what teachers already read in the classroom, and an ordinary webcam is enough to record it during online sessions. This made \ac{FER} an attractive option once deep learning became stable enough to be used outside the laboratory. Over the last decade, a dedicated line of educational \ac{FER} work has grown around engagement, attention, boredom, confusion, frustration, and related academic emotions, and a systematic review by Lek and Teo confirms that this line is now the dominant computer-vision approach to emotion monitoring in digital learning environments \cite{lek_academic_2023}. More recent syntheses also describe facial behaviour as one of the most practical signals for studying affect in classrooms and online courses \cite{vistorte_integrating_2024,florestiyanto_emotion_2024}.

A central piece of this development was the introduction of a public benchmark for in-the-wild educational affect. Gupta and colleagues released \ac{DAiSEE}, a dataset of 9{,}068 short webcam clips collected from 112 students, where each clip is labelled with four affective states, namely engagement, boredom, confusion and frustration, and each state is rated on four ordinal levels \cite{gupta_daisee_2016}. Their benchmark also reported the first deep-learning baselines on the dataset, using InceptionNet at frame and video level, a 3D convolutional network (C3D), and a long-term recurrent convolutional network (LRCN). Top-1 accuracy was clearly higher for confusion and frustration than for engagement and boredom, with LRCN reaching 53.7\% on boredom, 57.9\% on engagement, 72.3\% on confusion and 73.5\% on frustration, while frame- and video-level InceptionNet stayed below 50\% on boredom and engagement \cite{gupta_daisee_2016}. Two consequences shaped almost everything that came after. First, \ac{DAiSEE} became the reference benchmark for academic \ac{FER} \cite{lek_academic_2023}. Second, the dataset is genuinely difficult, especially for engagement and boredom, and is heavily class-imbalanced, which made temporal modelling, attention and imbalance-aware training the natural next steps. It is also worth noting that within the searched sources, no pure traditional machine-learning paper using \ac{DAiSEE} as the main method was found; classical components such as facial landmarks, action units or handcrafted features only appear inside hybrid pipelines combined with deep networks.

Early educational \ac{FER} work in the same period still kept one foot in classical computer vision. Rao and colleagues used a hybrid pipeline that merged Dlib landmark vectors with \ac{CNN} features and a support vector machine to classify the four \ac{DAiSEE} states, reporting 53.42\% overall accuracy on \ac{DAiSEE} alongside additional evaluation on JAFFE and CK+ \cite{rao_recognition_2021}. Around the same time, Leong compared FaceNet embeddings and landmark-distance streams, each followed by an \ac{LSTM}, to detect boredom and frustration on \ac{DAiSEE} \cite{leong_deep_2020}. These two studies showed that educational \ac{FER} had already moved away from generic basic-emotion labels toward states that matter for learning, but they also showed that single-frame or shallow-temporal architectures struggled with the hard \ac{DAiSEE} classes.

A clearer step forward came in 2021 with the work of Abedi and Khan, who combined a 2D ResNet-18 with a \ac{TCN} and trained it end-to-end on raw video frames. They followed the \ac{DAiSEE}-defined train, validation and test counts but merged training and validation for fitting, and reported 63.9\% four-class engagement accuracy on the test set, improving over a ResNet+\ac{LSTM} baseline at 61.15\% \cite{abedi_improving_2021}. Their study also made class imbalance an explicit design issue by using weighted cross-entropy and a balanced-batch sampling strategy. The same temporal direction was followed in educational \ac{FER} more broadly: Asaju and Vadapalli stacked a ResNet-152 with a \ac{BiLSTM} and mapped the four \ac{DAiSEE} states onto a simpler positive--neutral--negative pedagogical scheme, including a small live validation with real students \cite{asaju_adaptive_2022}. Together, these works confirmed that short video clips, rather than isolated frames, were a better unit of analysis for academic affect.

In 2022 the field moved toward attention and transformer-style processing of video. Mehta and colleagues introduced a 3D DenseNet built with spatial, temporal and hybrid non-local self-attention blocks and trained it with class-balanced focal and class-balanced MSE losses \cite{mehta_three-dimensional_2022}. They reported 63.59\% four-class engagement accuracy and 94.78\% binary-engagement accuracy on \ac{DAiSEE}, where the binary task is obtained by merging the two lowest engagement levels into ``low'' and the two highest into ``high'', and is therefore not directly comparable to the four-level results. They also reported a regression MSE of 0.0347 on \ac{DAiSEE} engagement and 0.0877 on EmotiW-EP. In the same year, Mandia et al. proposed a Vision Transformer with multi-head self-attention for automatic student engagement estimation on \ac{DAiSEE}, using imbalance-aware losses and reporting an overall accuracy of about 55.18\%, which the authors describe as outperforming Gupta et al.'s frame-level and video-level baselines by roughly 8--8.78 percentage points \cite{mandia_vision_2022}. The full PDF for this paper was not accessible, so its method details are limited to the available metadata and abstract.

By 2023, two related directions became visible. The first was the hybrid deep-classical pipeline, well represented by Santoni et al., who combined OpenFace-derived facial landmarks, action units, head-pose and gaze features, dimensionality reduction with PCA or SVD, oversampling with SMOTE, and a \ac{CNN} classifier \cite{santoni_convolutional_2023}. They reported a best maximum accuracy of 77.97\% with the SVD+\ac{CNN} configuration on a heavily resampled subset of \ac{DAiSEE} rather than on the original benchmark split, which means the number is not directly comparable with results that follow the standard \ac{DAiSEE} protocol. This work is important precisely because it shows that \ac{DAiSEE} research is not only about replacing classical methods with deep networks: it is also about combining deep features with classical preprocessing, feature engineering and imbalance handling. A complementary 2023 study, although not on \ac{DAiSEE}, illustrates the same hybrid trend in general \ac{FER}: Mukhiddinov et al.\ used MediaPipe landmarks together with \ac{HOG} features around those landmarks as input to a custom \ac{CNN} with batch normalisation and Mish activation, reporting 69.3\% accuracy across seven emotion classes on AffectNet under masked-face conditions \cite{mukhiddinov_masked_2023}. Both studies confirm that handcrafted facial cues remain useful when they are integrated as front-ends to deep models rather than used alone.

The 2024 work continued in three directions. Shiri et al.\ replaced earlier \ac{CNN} backbones with EfficientNetV2-L and combined it with \ac{GRU}, Bi-\ac{GRU}, \ac{LSTM} or Bi-\ac{LSTM} heads, reporting their best four-class engagement accuracy of 62.11\% on \ac{DAiSEE} with the \ac{LSTM} variant \cite{shiri_detection_2024}; the paper is sometimes cited under a co-author's name but the first author is Shiri. Tian et al.\ proposed a sequential ensemble model evaluated on both \ac{DAiSEE} and EmotiW-EP, reporting an MSE of 0.0386 on \ac{DAiSEE} and 0.0610 on EmotiW-EP, with the available abstract claiming stronger performance than prior approaches \cite{tian_predicting_2024}; the full PDF was not accessible, so this entry should be treated as limited verification. Su and colleagues then pushed the field further into the transformer era with a hybrid architecture, called DTransformer, that combines a part-and-sensitive attention network on ResNet-18 patches with a 3D-\ac{CNN} tube tokenizer feeding a multi-head self-attention transformer, fusing both branches for prediction \cite{su_leveraging_2024}. They reported 64.0\% four-class engagement accuracy on \ac{DAiSEE} and a regression MSE of 0.0729 on EmotiW-EP, which makes their work one of the strongest current transformer-based engagement detectors.

Two recent 2025 studies show that hybrid deep-classical thinking is still alive. Rianto et al.\ took a graph-based view: they extracted 68 facial landmarks per frame, built a Delaunay triangulation graph, fed it to a graph convolutional network, and combined it with a 1D \ac{CNN} over 52 action-unit values, training the bimodal model with the LDAM imbalance-aware loss \cite{rianto_engagement_2025}. They followed the \ac{DAiSEE} engagement-only split and reported a best accuracy of 52.24\% with macro-F1 of 0.49, while honestly noting that minority engagement levels remained largely mispredicted. In general \ac{FER}, Türkeş and Aydın combined a pre-trained ResNet-50 feature extractor with a Support Vector Machine classifier \cite{turkes_enhanced_2025}. The \ac{DAiSEE} branch is therefore now bracketed by hybrid graph-and-\ac{CNN} methods on one side, and by classical hybrid \ac{CNN}+\ac{SVM} pipelines on standard \ac{FER} datasets on the other.

Educational \ac{FER} has also matured in its target applications. Komaravalli and Janet fine-tuned an Xception-based transfer-learning model on a rebalanced \ac{DAiSEE} subset and cross-tested it on OL-SFED \cite{komaravalli_detecting_2023}. Gupta and colleagues proposed EDFA, an ensemble that fuses VGG19 and ResNet50 outputs into a rule-based attention index for real-time classroom monitoring, deployed as a web application \cite{gupta_edfa_2023}. Aly combined a ResNet-50 backbone with the Convolutional Block Attention Module and a Temporal Convolutional Network for online student monitoring, reporting throughput of about 28 frames per second on a GPU, 12 on a CPU, and 15 on a Jetson Nano \cite{aly_revolutionizing_2025}, and a follow-up study added a 3D \ac{CNN} branch with an optimised fusion layer \cite{aly_comprehensive_2025}. Tulsani and colleagues used a hybrid of \ac{MTCNN} face detection, a Vision Transformer and a \ac{BiLSTM} inside a learning-management system, with a near-real-time latency of about 0.8 seconds per frame sequence \cite{tulsani_realtime_2025}. Liu and colleagues proposed an adaptive global attention block together with a differential temporal transformer, evaluated on \ac{DAiSEE}, DFEW and FERV39k, and reported about 11.7 milliseconds per GPU inference \cite{liu_dynamic_2026}. Woolf and colleagues described the Face Readers line of work inside the MathSpring intelligent tutoring system, where attention, head pose and facial cues are used to predict problem-solving outcomes and to trigger real-time interventions \cite{woolf_face_2023}. Together, these examples show that educational \ac{FER} is no longer only about classification accuracy; it is slowly being connected to the feedback loops of adaptive learning systems.

Taken together, the \ac{DAiSEE} evidence base evolved through deep-learning baselines, \ac{CNN}-plus-temporal models, 3D and self-attention architectures, hybrid deep-classical pipelines, EfficientNet plus RNN systems, transformer-based detectors and ensemble models, and was complemented by hybrid \ac{CNN}-and-classical \ac{FER} work on standard expression datasets. Despite this progress, several limits remain. Datasets for academic emotions are still small and imbalanced; results on the four-level engagement task remain in the low-to-mid sixties even for strong transformer or attention models, and many studies do not use the official \ac{DAiSEE} split or do not report per-class metrics, which makes head-to-head comparison difficult \cite{abedi_improving_2021,santoni_convolutional_2023,rianto_engagement_2025,shiri_detection_2024,su_leveraging_2024}. Generalisation across cultures and recording conditions remains weak, and visually similar states such as confusion and frustration are still confused \cite{lek_academic_2023,florestiyanto_emotion_2024,vistorte_integrating_2024}. More importantly for this thesis, almost all of the studies above stop at engagement classification, regression or visualisation; very few connect the \ac{FER} output to a concrete next-task selection or adaptation policy, and live deployment in real classrooms is reported only in a handful of cases. Despite these concerns, facial behaviour remains one of the few signals that can be captured continuously, without attaching sensors to students, without interrupting the lesson, and with hardware that most classrooms and online learners already have. This practical profile is one of the reasons why \ac{FER}, on its own, continues to be central to recent educational-affect research, and why a webcam-only \ac{FER} pipeline that is explicitly connected to an adaptive next-task selector still defines an open and useful design point. Table~\ref{tab:fer_comparison} summarises representative \ac{DAiSEE}-based \ac{FER} studies for benchmark comparison. The table spans two landscape pages: the first lists each study, dataset, and task; the second gives the full architecture description and evaluation metrics for the same studies.



\clearpage
\newgeometry{left=3.2cm,right=2.5cm,top=2.5cm,bottom=2.5cm}
% --- Table~\ref{tab:fer_comparison} part 1: study descriptors (landscape page 1) ---
\begin{landscape}
\thispagestyle{empty}
\vspace*{-0.5cm}
\large
\setlength{\tabcolsep}{5pt}
\renewcommand{\arraystretch}{1.2}
\setlength{\LTleft}{0pt}
\setlength{\LTright}{0pt}
\begin{longtable}{>{\RaggedRight\arraybackslash}p{0.16\linewidth}
                  >{\RaggedRight\arraybackslash}p{0.14\linewidth}
                  >{\RaggedRight\arraybackslash}p{0.66\linewidth}}
\caption{Comparison of representative \ac{DAiSEE}-based \ac{FER} studies}
\label{tab:fer_comparison} \\

\toprule
\textbf{Study} & \textbf{Dataset} & \textbf{Task} \\
\midrule
\endfirsthead

\toprule
\textbf{Study} & \textbf{Dataset} & \textbf{Task} \\
\midrule
\endhead

\bottomrule
\endfoot

\makecell[l]{Gupta et al.\\ (2016) \cite{gupta_daisee_2016}} &
\ac{DAiSEE} &
4 academic states, 4 ordinal levels \\

\makecell[l]{Abedi and Khan\\ (2021) \cite{abedi_improving_2021}} &
\ac{DAiSEE} &
4-class engagement classification \\

\makecell[l]{Mehta et al.\\ (2022) \cite{mehta_three-dimensional_2022}} &
\ac{DAiSEE}, EmotiW-EP &
4 academic states + engagement regression \\

\makecell[l]{Su, He, Luo\\ (2024) \cite{su_leveraging_2024}} &
\ac{DAiSEE}, EmotiW-EP &
4-class engagement (\ac{DAiSEE}); regression (EmotiW-EP) \\

\makecell[l]{Liu et al.\\ (2026) \cite{liu_dynamic_2026}} &
\ac{DAiSEE}, DFEW, FERV39k &
4 academic states (in-the-wild video \ac{FER}) \\

\makecell[l]{Rianto et al.\\ (2025) \cite{rianto_engagement_2025}} &
\ac{DAiSEE} &
4-class engagement; imbalance-focused \\

\end{longtable}
\end{landscape}

% --- Table~\ref{tab:fer_comparison} part 2: architecture and metrics (landscape page 2) ---
\clearpage
\begin{landscape}
\thispagestyle{empty}
\vspace*{-0.5cm}
\large
\setlength{\tabcolsep}{5pt}
\renewcommand{\arraystretch}{1.2}
\setlength{\LTleft}{0pt}
\setlength{\LTright}{0pt}
\addtocounter{table}{-1}
\begin{longtable}{>{\RaggedRight\arraybackslash}p{0.14\linewidth}
                  >{\RaggedRight\arraybackslash}p{0.58\linewidth}
                  >{\RaggedRight\arraybackslash}p{0.24\linewidth}}
\multicolumn{3}{l}{\textit{Table~\ref{tab:fer_comparison} (continued): architectural details and evaluation metrics.}} \\[0.4em]

\toprule
\textbf{Study} & \textbf{Architecture} & \textbf{Evaluation Metric(s)} \\
\midrule
\endfirsthead

\multicolumn{3}{l}{\textit{Table~\ref{tab:fer_comparison} (continued).}} \\[0.4em]
\toprule
\textbf{Study} & \textbf{Architecture} & \textbf{Evaluation Metric(s)} \\
\midrule
\endhead

\bottomrule
\multicolumn{3}{p{0.98\linewidth}}{\textit{Note.} Each architecture entry summarises the full pipeline as reported by the original study. Metrics are reported as published; four-class engagement accuracy, per-class accuracy, and regression MSE are not directly comparable across rows. Abbreviations: B, boredom; E, engagement; C, confusion; F, frustration.} \\
\endfoot

\makecell[l]{Gupta et al.\\ (2016)} &
Backbone: InceptionNet (frame/video), C3D (3D \ac{CNN}), LRCN (\ac{CNN} encoder).\newline
Temporal: 3D convolution (C3D) or \ac{CNN}+\ac{LSTM} (LRCN).\newline
Attention: none.\newline
Classification: per-state ordinal top-1 (4 labels). &
Per-class top-1 acc.\ (LRCN): B 53.7\%, E 57.9\%, C 72.3\%, F 73.5\% \\

\makecell[l]{Abedi and Khan\\ (2021)} &
Backbone: ResNet-18 (ImageNet-pretrained).\newline
Temporal: dilated \ac{TCN} over frame features.\newline
Attention: none.\newline
Classification: 4-class softmax; weighted cross-entropy and balanced-batch sampling. &
63.9\% 4-class acc. \\

\makecell[l]{Mehta et al.\\ (2022)} &
Backbone: 3D DenseNet-121.\newline
Temporal: 3D convolution (tube modelling).\newline
Attention: spatial, temporal, and hybrid non-local self-attention blocks.\newline
Classification: 4-class + regression heads; class-balanced focal and CB-MSE losses. &
63.59\% 4-class acc.; MSE 0.0347 (regression) \\

\makecell[l]{Su, He, Luo\\ (2024)} &
Backbone: ResNet-18 patches (PANet) + 3D \ac{CNN} tube tokenizer (STformer).\newline
Temporal: 3D \ac{CNN} tubes + transformer encoder.\newline
Attention: part-sensitive spatial attention + multi-head self-attention; late fusion.\newline
Classification: 4-class engagement / regression; dual-branch fusion head. &
64.0\% 4-class acc.\ (\ac{DAiSEE}); MSE 0.0729 \\

\makecell[l]{Liu et al.\\ (2026)} &
Backbone: ResNet-18.\newline
Temporal: Differential Temporal Transformer (DTFormer).\newline
Attention: Adaptive Global Attention block.\newline
Classification: 4-class softmax over fused spatiotemporal features. &
62.56\% (\ac{DAiSEE}, 4-class); evaluated also on DFEW and FERV39k (cross-dataset \ac{FER}) \\

\makecell[l]{Rianto et al.\\ (2025)} &
Backbone: Delaunay GCN on 68 landmarks + 1D-\ac{CNN} on 52 action units (bimodal).\newline
Temporal: none (per-frame streams).\newline
Attention: none.\newline
Classification: 4-class engagement; LDAM imbalance-aware loss. &
52.24\% acc.; macro-F1 0.49 \\

\end{longtable}
\end{landscape}
\restoregeometry
\clearpage

\section{Adaptive Learning Systems}\label{sec:adaptive_learning_systems}

While the previous section focused on recognising learner states through facial behaviour, adaptive learning systems focus on deciding how instruction should respond to learner needs. These systems are digital tools that change how they teach depending on the individual learner. Instead of presenting the same material in the same order to every student, they adjust things like the content, the pace, the difficulty of the tasks, the feedback given after each answer, hints, and presentation format \cite{strielkowski_span_2025}. A common way to describe this behaviour is as a closed loop: the system collects data from the learner, uses that data to judge how the learner is progressing, and then decides what to show next \cite{gligorea_adaptive_2023}. Reviews of intelligent tutoring systems stress that adaptation is not limited to picking the next exercise: the tutor model may also schedule scaffolding, remediation, motivational messages, interface changes, and assessment pacing. The goal is to give each learner tasks that match their current level, so that they can keep improving without feeling bored or overwhelmed \cite{zerkouk_comprehensive_2025}.

Early adaptive systems focused on explicit rules written by experts. These were called \ac{ITS}, and examples such as GUIDON and the LISP Tutor were built in the 1980s around hand-crafted if--then rules that decided when to give a hint, when to move on, and when to mark a concept as learned . The strength of these systems was that their decisions were transparent and easy to explain, but the rules had to be prepared for every topic, which made them hard to extend beyond a single subject. Rule-based logic has not disappeared, though. Sayed et al., for example, used a small set of threshold rules based on a student's penalised grade to move an exercise up or down in difficulty inside an online platform for third-grade mathematics \cite{sayed_ai-based_2023}.

Later work introduced student models that tried to estimate what each learner already knew. Instead of relying only on teacher-written rules, these systems built a picture of the learner from their interaction history and updated it after every answer \cite{zerkouk_comprehensive_2025}. This line of research eventually led to knowledge tracing methods, in which the system predicts how likely the learner is to answer the next question correctly based on their earlier answers. Pardos and colleagues showed that year-long patterns of boredom, engaged concentration, confusion, and gaming-the-system behaviour in ASSISTments logs correlate with end-of-year mathematics test scores \cite{Pardos2014AffectiveStatesStateTests}. Kuling and Zitnik built a recent exercise recommender that tries to pick questions just above the learner's current level \cite{kuling_ken_2025}. Educational psychologists call this range the Zone of Proximal Development: the set of tasks that are challenging but still reachable. Their system ranks candidate questions by predicted difficulty and then chooses items around the 66th percentile, so that the next exercise is hard but not impossible \cite{kuling_ken_2025}.

Over time, researchers began to treat task selection as a decision problem that could be learned from experience. \ac{RL} then became a common tool for this. In \ac{RL}, the system tries different actions, observes the results, and slowly learns which actions lead to better outcomes. Azoulay et al. compared several adaptation methods within the same simulated environment. The study tested different reinforcement learning algorithms, a bandit-style method, and a Bayesian approach that keeps a probability estimate of each learner's ability \cite{azoulay_adaptive_2021}. The Bayesian method reached about 89 to 93 percent of the performance of an ideal algorithm that already knew each student's real ability, while the reinforcement learning methods reached lower but still useful levels \cite{azoulay_adaptive_2021}. The study has become an important reference because it placed very different adaptation families on equal ground, and later adaptive tutoring work extended these baselines \cite{kubotani_rltutor_2021,fu_integrating_2025}.

After Azoulay et al.\ showed in 2021 that tutoring adaptation could be studied as a sequential decision problem inside a shared simulator, the next step was to ask whether a policy could be trained against an explicit learner model rather than only compared as a fixed algorithm \cite{azoulay_adaptive_2021,kubotani_rltutor_2021}. Kubotani et al.\ (2021) addressed spaced review scheduling, where the tutor must decide which previously studied item to present before forgetting occurs. They replaced live learners with a DAS3H-based mathematical student model and trained a \ac{PPO} agent to select review items over 30 items across 15 days and 30 sessions. The state came from estimated memory rates produced by the inner student model, and the action space selected the next item to review. Evaluation used only the mathematical simulator and compared RLTutor against random, Leitner, and threshold teachers; RLTutor maintained higher average retention in early sessions, although the gap narrowed later. The study matters because it shows why synthetic learners appear before classroom deployment, but its student model was still narrow and its action space was limited to review selection rather than full tutoring control.

The same year, Jezuina Koroveshi and Ktona (2021) moved closer to the architecture used in the present thesis by training a pedagogical module with a Deep Q-Network in a Python intelligent tutoring system \cite{jezuina_koroveshi_training_2021}. Their educational problem was lesson sequencing: the tutor had to decide whether a learner should remain in the current lesson or advance after concept mastery. The state was a binary vector encoding the active lesson, the concepts taught in that lesson, and the concepts already learned by the simulated student. The student model was a simple simulated learner with a fixed learning probability of 0.7, where each taught concept was learned or missed according to that probability. The action space contained two choices: stay in the current lesson or move to the next lesson. The reward function gave $+10$ when the action matched the mastery condition and $-10$ when it did not. Training used \ac{DQN} with experience replay, a target network, batch size 64, discount factor $\gamma=0.85$, and epsilon decay over up to 4{,}000 episodes in two scenarios, one with empty prior knowledge and one with random prior knowledge. Evaluation reported episode reward and episode length curves over 100 test episodes, but did not compare against a rule-based tutor baseline. This paper is the earliest clear \ac{DQN} tutoring design in the set reviewed here, although its state did not include knowledge tracing or engagement sensing.

By 2023, the question shifted from whether \ac{DQN} could train in simulation to whether any part of the tutoring policy could survive contact with a real classroom. Sayed et al.\ (2023) built the APPEAL platform for grade-three mathematics and split adaptation into two decision layers \cite{sayed_ai-based_2023}. Exercise difficulty and navigation were controlled by an online rule-based algorithm using penalised grades, attempts, and hints across Bloom taxonomy levels. Content presentation was handled separately by a Deep Q-Network formulated as a Markov decision process over student interaction traces. The \ac{DQN} state included time spent per lesson, exercise scores, hint use, and repeated attempts within taxonomy levels and difficulty bands. The action space encoded VARK presentation modes as four binary switches that showed or hid visual, aural, read/write, and kinesthetic material. The reward combined academic improvement with savings in time, hints, and attempts. Training used \ac{DQN} with experience replay; live evaluation ran for two months with 26 learners and reported post-test gains relative to pre-test, with the largest improvements for students who started with lower performance. The study shows that \ac{DQN} can operate inside a deployed school platform, but only for presentation choice, not for full tutoring control, and the sample remained small.

Once \ac{DQN} tutors existed both in simulation and in partial live deployment, the next development in 2025 was to strengthen the student model before policy learning. Fu (2025) proposed \ac{RL}-DKT, which uses a recurrent dynamic knowledge tracing network to estimate skill mastery and a REINFORCE agent to select the next task from offline logs \cite{fu_integrating_2025}. The educational problem was learning-path sequencing rather than live hint or feedback control. The state was the DKT hidden vector updated after each response; the action space was the set of candidate tasks available at each step; and the reward combined correctness, task difficulty, and progress terms rather than a single binary outcome. Evaluation used ASSISTments, KDD Cup 2010, and Cognitive Tutor logs and reported AUC of 0.874, F1 of 0.812, average completion time of 34.5 minutes, dropout of 3.2\%, and a reported 12.5\% reduction in completion time relative to baselines. \ac{RL}-DKT is not a \ac{DQN} tutor, but it clarified an important design pattern for this thesis: separate student modelling from policy optimisation. Its limitation is offline replay rather than an online tutor loop with synthetic or live students.

Guzman and Cruz-Mercado (2025) extended that pattern with a larger action space and a two-stage evaluation design \cite{guzman_developing_2025}. Their \ac{RL}-Tutor system targeted introductory Python skill development. A Dynamic Key-Value Memory Network maintained the student model from correctness, response time, and attempts, producing a knowledge-tracing vector concatenated with recent learning-object context and session meta-features such as consecutive errors and hint requests. A \ac{PPO} agent then selected hierarchical actions: problem versus instructional page versus review quiz, difficulty level, target skill, hint type, and pacing decision. The reward combined performance, efficiency, learning gain, engagement, and retention terms. Because online \ac{RL} sample collection is slow, the agent was first trained for about 400{,}000 steps against 10{,}000 synthetic students and then tested in a between-subjects study with 90 human learners against rule-based and static tutors. Simulated mean cumulative reward reached $315 \pm 22$ versus $205 \pm 15$ for the rule-based tutor ($p<0.001$). Live normalised learning gain was 0.72 for \ac{RL}-Tutor, 0.58 for the rule-based tutor, and 0.45 for the static tutor ($p<0.01$), with similar retention advantages one week later. This study improved both the student model and the action space, but it used \ac{PPO} rather than \ac{DQN} and still inferred engagement only from behaviour, not from facial signals.

The most recent \ac{DQN} work in this line returns to large-scale synthetic evaluation. Zou and Xiong (2026) built a Double \ac{DQN} path optimiser for a knowledge-graph curriculum with 537 nodes and 1{,}243 relations \cite{zou_deep_2026}. The state represented learner mastery and progress over the graph; the action space selected the next course item; and the reward combined learning effectiveness, time cost, and knowledge coverage. Training used experience replay and epsilon-greedy exploration in a simulation with 5{,}000 virtual learners. The system reached an 87.3\% completion rate, mastery of 0.823, and outperformed random, collaborative-filtering, and standard Q-learning baselines. For beginner learners, average mastery rose from 0.512 to 0.789 in ablation analysis. This paper is the strongest recent example of \ac{DQN} plus synthetic cohort evaluation, but it optimises path order rather than hints, feedback, or pacing, and it was not followed by a classroom study.

Viewed across these dates, the \ac{DQN} tutoring line changed in a clear order. Jezuina Koroveshi and Ktona (2021) first mapped lesson progression to \ac{DQN} actions in a simulated Python tutor with a simple mastery vector and a fixed reward table \cite{jezuina_koroveshi_training_2021}. Sayed et al.\ (2023) kept \ac{DQN} for presentation only and left difficulty to rules, which reduced the learned portion of the policy but provided rare live-school evidence \cite{sayed_ai-based_2023}. Zou and Xiong (2026) then moved to Double \ac{DQN} with a structured knowledge-graph state and 5{,}000 virtual learners, improving scale but narrowing the problem to path selection \cite{zou_deep_2026}. What improved over time was the realism of evaluation settings and the richness of student models in hybrid systems such as \ac{RL}-DKT and \ac{RL}-Tutor \cite{fu_integrating_2025,guzman_developing_2025}. What remained unresolved was a \ac{DQN} tutor that combines an explicit student model, a multi-action pedagogical policy, and reusable synthetic evaluation without splitting control across rules and learned components.

\subsection{Simulation, Synthetic Students, and the Ethics of Live Policy Learning}\label{subsec:synthetic_student_evaluation}

Deep reinforcement-learning algorithms require interaction volumes that real classrooms structurally cannot supply during policy development. Sanz Ausin and colleagues confronted this directly: their offline \ac{ITS} corpus contained fewer than 800 complete student trajectories, a fraction of the interaction budget that \ac{DQN} typically consumes on even simple benchmarks \cite{sanz_ausin_leveraging_2019}. Dulac-Arnold and colleagues formalise the same constraint for production systems more broadly, arguing that online policy search must be both sample-efficient and rapidly performant because exploratory actions on real users carry immediate cost \cite{dulacarnold_challenges_2021}. Kubotani et al.\ reach the same conclusion for adaptive review scheduling: the interaction count required for \ac{PPO} optimisation cannot be extracted from live learners without prohibitive delay, which is precisely why RLTutor trains against a virtual student model \cite{kubotani_rltutor_2021}. Li's survey of \ac{RL} in education confirms that sample intensity and exploration risk remain persistent barriers across algorithmic families \cite{li_reinforcement_2022}. The gap between what policy search demands and what a classroom can supply is structural, not incidental.

The practical constraints extend beyond data volume. During policy search, a tutor must try actions whose value is still uncertain, and delayed, noisy educational outcomes make harmful policies difficult to detect after deployment \cite{dulacarnold_challenges_2021}. Reward misspecification compounds this risk: a proxy that over-weights engagement can produce policies that reach high measured reward while delivering negligible mastery gain. Olukola and Rahimi demonstrate this concretely in a simulated Python tutor, where an engagement-heavy policy drove the Reward Hacking Severity Index to 0.317; architectural safety constraints reduced it to 0.102 \cite{olukola_pedagogical_2026}. Exploratory policy learning on real learners is therefore constrained by safety, cost, and reward uncertainty, not by a categorical prohibition on classroom research. The important distinction is between the training phase, where a policy is still searching and reward signals are uncertain, and the evaluation phase, where a fixed, already-validated policy is tested under ethical approval. Sanz Ausin et al.\ follow exactly this separation: offline deep \ac{RL} on historical logs, followed by randomised classroom studies once a policy is selected \cite{sanz_ausin_leveraging_2019}.

Because of these constraints, synthetic student environments have become the standard first-stage training platform in \ac{RL} tutoring research, and the models used to populate them have grown progressively richer. Early \ac{DQN} work by Jezuina Koroveshi and Ktona trained a lesson-progression policy against a simulated learner with a fixed concept-learning probability across thousands of simulated episodes \cite{jezuina_koroveshi_training_2021}. Kubotani et al.\ replaced that with a DAS3H memory model pre-trained on EdNet, so the \ac{RL} agent optimised review scheduling by interacting with a calibrated inner student model rather than with students directly \cite{kubotani_rltutor_2021}. Guzman and Cruz-Mercado extended this design with a DKVMN synthetic cohort of 10{,}000 virtual learners and approximately 400{,}000 \ac{PPO} training steps before a between-subjects study with 90 human participants \cite{guzman_developing_2025}. At larger scale, Azoulay et al.\ used a shared virtual-student environment to benchmark competing adaptation families under identical conditions \cite{azoulay_adaptive_2021}, and Zou and Xiong sampled 5{,}000 virtual learners to evaluate a Double \ac{DQN} path optimiser on a 537-node knowledge graph \cite{zou_deep_2026}. Across all of these systems, simulation provides the controllable, reproducible, and scalable stage that precedes any contact with real classrooms and from which policies graduate to live evaluation.

The present thesis follows this simulation-first pipeline. Existing synthetic environments represent learners through knowledge-tracing estimates \cite{azoulay_adaptive_2021,fu_integrating_2025,kubotani_rltutor_2021,jezuina_koroveshi_training_2021}; the next natural step is to extend the synthetic student state to include a multidimensional affective state representation derived from facial expression recognition. Rather than treating emotion as a single engagement proxy, the thesis integrates four discrete academic emotions identified in prior educational \ac{FER} work: boredom, confusion, engagement, and frustration. This affective state representation is integrated into the \ac{DQN} state space and trained inside a synthetic student environment before any classroom use. Table~\ref{tab:adaptive_systems_comparison} summarises the studies that most directly inform this design.

\clearpage
\newgeometry{left=1.5cm,right=1.5cm,top=2.0cm,bottom=2.0cm}
\begin{landscape}
\thispagestyle{empty}
\setlength{\tabcolsep}{5pt}
\renewcommand{\arraystretch}{1.35}
\small

\begin{table}[h]
\caption{Simulation-based \ac{RL} tutoring studies most relevant to the present thesis design}
\label{tab:adaptive_systems_comparison}
\begin{tabularx}{\linewidth}{%
  >{\RaggedRight\arraybackslash}p{2.4cm}
  >{\RaggedRight\arraybackslash}p{4.0cm}
  >{\RaggedRight\arraybackslash}p{2.2cm}
  >{\RaggedRight\arraybackslash}X
  >{\RaggedRight\arraybackslash}p{3.8cm}
  >{\RaggedRight\arraybackslash}p{4.2cm}}
\toprule
\textbf{Paper} & \textbf{Student model} & \textbf{\ac{RL} method} & \textbf{Learning task} & \textbf{Evaluation setup} & \textbf{Main result} \\
\midrule

\makecell[l]{Jezuina Koroveshi\\ \& Ktona (2021)\\ \cite{jezuina_koroveshi_training_2021}} &
A simple simulated learner that learns with a fixed probability &
\ac{DQN} &
Deciding whether to stay in the current lesson or advance (Python) &
Simulation only &
Learning improved steadily in simulation training; no live student test \\[4pt]

\makecell[l]{Zou \& Xiong\\ (2026)\\ \cite{zou_deep_2026}} &
5{,}000 simulated learners across different ability levels &
\ac{DQN} &
Selecting the next course item from a knowledge graph &
Large-scale simulation only &
87.3\% course completion; strong learning performance (mastery 0.823) \\[4pt]

\makecell[l]{Sanz Ausin\\ et al.\ (2019)\\ \cite{sanz_ausin_leveraging_2019}} &
Trained on historical student data (no live students used) &
\ac{DQN} (offline) &
Choosing worked examples vs.\ problem-solving in logic &
Trained offline, then tested with real students &
Significant learning gain for stronger learners; real-time training with students is not practical \\[4pt]

\makecell[l]{Guzman \&\\ Cruz-Mercado\\ (2025)\\ \cite{guzman_developing_2025}} &
A knowledge-tracing model trained on 10{,}000 simulated learners &
Reinforcement learning &
Selecting problem type, difficulty, hints, and pacing (Python) &
Trained in simulation, then tested with 90 real students &
Better learning quality (0.72 vs.\ 0.58) than a rule-based tutor \\[4pt]

\makecell[l]{Sayed\\ et al.\ (2023)\\ \cite{sayed_ai-based_2023}} &
Real classroom learners (no simulated cohort) &
\ac{DQN} &
Choosing how content is presented (visual, audio, read/write, hands-on) &
Two-month classroom pilot with 26 students &
Post-test gains, largest for lower-performing students \\

\bottomrule
\multicolumn{6}{p{\linewidth}}{\textit{Note.} All studies train \ac{RL} tutoring policies using simulated students or historical data, since exploratory \ac{RL} training is typically performed in simulation or offline due to safety and cost constraints. Real-student studies are used only for controlled evaluation after a policy has been validated. None of these systems uses facial expression recognition in the learner state.}
\end{tabularx}
\end{table}

\end{landscape}
\restoregeometry
\clearpage
\section{Emotion-Aware Adaptation}\label{sec:emotion_aware_adaptation}

The previous section reviewed reinforcement-learning tutors that adapt from knowledge and performance traces alone. This section looks at systems that also use learner emotion. The work in this area has shifted over time. Early systems read emotion from costly multi-sensor rigs or from weak text signals. Later systems moved toward webcam \ac{FER} joined with learned adaptation policies. That shift is where the present thesis sits: a single webcam reads the face, and a learned policy decides what to do next. The studies below trace that shift and report what each system actually measured.

The first practical systems read emotion from dialogue with the tutor. AutoTutor held short natural-language dialogues and asked whether dialogue signals could show how the learner felt \cite{Graesser2004AutoTutorDialogue}. In 2005, D'Mello and colleagues combined affect sensors with AutoTutor to identify boredom, confusion, frustration and engagement \cite{dmello_integrating_2005}. Later work showed that ordinary features of tutoring conversation carried information about boredom, confusion, flow and frustration \cite{DMello2008AutomaticDetectionLearnersAffect,dmello_language_2012}, and that confusion, when handled well, can help learning \cite{dmello2014confusionbeneficial}. This was the historical start of the field: emotion can be detected during instruction. The response, at that stage, was small hint and encouragement changes. A more sensor-heavy line followed. Woolf and colleagues built the Wayang Outpost tutor with a four-sensor rig: a webcam, a pressure-sensitive chair, a pressure mouse and a skin-conductance bracelet \cite{woolf_affect-aware_2009}. It estimated states such as frustration and engagement and changed problems, encouragement, or an animated learning companion. The reported direction was less frustration and more interest when companions were used, and the effect was stronger for girls. This shows that multi-sensor emotion-aware tutoring is old, but it was costly: four sensors per student, custom hardware, a sensor server and a research team for one school. A cheaper line then tried a webcam with text. Lin and colleagues combined \ac{FER} with text analysis in a tutor for digital art, using an Active Shape Model on webcam images and a semantic voting method on typed input \cite{lin_employing_2012}; the evaluation tested usability only, not learning gains. Fwa proposed a rule-based affective tutor for novice programmers that mapped detected emotions to predefined tutoring strategies \cite{fwa_architectural_2018}. A classroom trial with 39 learners (21 control, 18 test) reported faster task completion (median 470.5 s versus 831 s, $p<0.05$) and more exercises attempted (median 7.5 versus 5, $p<0.05$). Leong then studied efficient deep models that detect boredom and frustration from webcam video in real time, which made real-time webcam reading realistic without extra hardware \cite{leong_deep_2020}.

Petrovica and Ekenel reviewed fourteen affective tutoring systems and named the trade-off that still shapes the field \cite{petrovica_emotion_2016}. They described three paths. Sensor-free log methods scale but lose accuracy. Sensor-heavy physiological or camera methods are accurate but costly and uncomfortable. They proposed a sensor-lite middle path: use only the webcam or microphone that classrooms already have. They also reported that most systems of that period lacked rich real-time adaptation.

Two systems around 2020 sat at the two ends of this trade-off. Megahed and Mohammed combined \ac{FER} with fuzzy-logic rules to change difficulty and presentation \cite{megahed_modeling_2020}; no learning result was reported. Standen and colleagues evaluated the MaTHiSiS system for learners with intellectual disabilities, using five sensing channels (face, gaze, body pose, voice and a gesture accelerometer) \cite{standen_evaluation_2020}. The study ran with 67 learners across six schools in the United Kingdom, Italy and Spain. Sessions driven by affect and achievement gave significantly more engagement and significantly less boredom than achievement-only sessions, but there was no significant difference in learning gains.

From 2024 onward, several systems connected emotion recognition to reinforcement learning \cite{perez_emotions_2024,govea_implementation_2024,gong_design_2025,kong_real-time_2025}. Pérez, Dapena and Aguilar used two Q-learning tables over a circumplex model and chose when to exploit or explore based on the learner's emotional region \cite{perez_emotions_2024}. The affect data came from a static expression corpus, not from a live webcam module. The method was tested in simulation only, with one thousand runs over two hundred tasks and no human participants, and reached a learning gain of 0.49 at task 200, above all baselines (next best 0.32). Govea and colleagues used multimodal sensing (cameras, microphones and biometric sensors) with deep reinforcement learning, trained on data from 500 students; the reported emotion detection accuracy rose to 89.3\%, with reported gains in real-time adaptability and personalisation \cite{govea_implementation_2024}. Kong and colleagues used multimodal sensing (webcam, voice, physiological signals and logs) with a reinforcement-learning controller built as a \ac{DQN}-backed contextual bandit \cite{kong_real-time_2025}. A lab and real-classroom randomised study with 271 learners reported affect recognition accuracy of 82 to 88\% and about 15 percentage points higher learning gain for mid-ability learners over a static baseline.

Gong and colleagues built a \ac{DQN} multimodal system (webcam, microphone and text) that detects basic emotions, namely happy, sad, angry, neutral and surprised, rather than academic emotions \cite{gong_design_2025}. It reported a fusion accuracy of 0.91 and an AUC of 0.94. These numbers come with clear limits: the system was evaluated on public datasets with no real students, and it targets basic emotions rather than the academic emotions that matter for learning. A scoping review of twenty-seven systems found that facial expression was the dominant sensing channel and scaffolding the most common pedagogical response \cite{fernandez-herrero_evaluating_2024}.

Table~\ref{tab:affective_systems_comparison} compares these systems and points to the gap addressed in the next section: emotion is rarely joined to a learned, webcam-only tutoring policy.

\clearpage
\newgeometry{left=1.5cm,right=1.5cm,top=2.0cm,bottom=2.0cm}
\begin{landscape}
\thispagestyle{empty}
\small
\setlength{\tabcolsep}{4pt}
\renewcommand{\arraystretch}{1.5}
\setlength{\LTleft}{0pt}\setlength{\LTright}{0pt}

\begin{longtable}{>{\RaggedRight\arraybackslash}p{0.10\linewidth}
                  >{\RaggedRight\arraybackslash}p{0.17\linewidth}
                  >{\RaggedRight\arraybackslash}p{0.18\linewidth}
                  >{\RaggedRight\arraybackslash}p{0.15\linewidth}
                  >{\RaggedRight\arraybackslash}p{0.15\linewidth}
                  >{\RaggedRight\arraybackslash}p{0.21\linewidth}}
\caption{Emotion-aware tutoring systems: sensing, adaptation, and reported results.}
\label{tab:affective_systems_comparison} \\

\toprule
\textbf{Study} & \textbf{Sensing modality} & \textbf{Emotions detected} & \textbf{Adaptation method} & \textbf{Deployment setting} & \textbf{Main result} \\
\midrule
\endfirsthead

\toprule
\textbf{Study} & \textbf{Sensing modality} & \textbf{Emotions detected} & \textbf{Adaptation method} & \textbf{Deployment setting} & \textbf{Main result} \\
\midrule
\endhead

\bottomrule
\multicolumn{6}{p{0.98\linewidth}}{\small\textit{Note.} Facial expression is the most common sensing channel across these systems. Several systems add extra hardware cost. A webcam-only \ac{FER} signal joined to a learned policy is the configuration this thesis targets.} \\
\endfoot



\makecell[l]{Fwa\\ (2018)\\ \cite{fwa_architectural_2018}}
& webcam plus keyboard/mouse logs plus eye tracker
& frustration, (dis)engagement
& rule-based (Rete engine)
& real classroom trial, n=39
& faster task completion (median 470.5 s vs 831 s, p<0.05); more exercises attempted \\


\makecell[l]{Standen\\ (2020)\\ \cite{standen_evaluation_2020}}
& five sensors: face, gaze, body pose, voice, gesture
& engagement, boredom
& affect plus achievement rules
& six schools, 67 learners
& more engagement and less boredom vs achievement-only; no learning gain difference \\

\makecell[l]{P\'erez\\ (2024)\\ \cite{perez_emotions_2024}}
& static expression corpus (not a live webcam)
& valence, arousal (anxiety, boredom)
& Q-learning, two tables
& simulation only (no students)
& learning gain 0.49 at task 200, above all baselines (next best 0.32) \\

\makecell[l]{Govea\\ (2024)\\ \cite{govea_implementation_2024}}
& multimodal: cameras, microphones, biometric sensors
& happiness, sadness, anger, surprise, fear, disgust (six basic)
& deep \ac{RL}
& data from 500 students
& emotion detection accuracy rose from 72.4 to 89.3\% (k-fold, k=10) \\

\makecell[l]{Kong\\ (2025)\\ \cite{kong_real-time_2025}}
& multimodal: webcam, voice, physiological, logs
& engagement, confusion, frustration, boredom, delight, anxiety, neutral, surprise
& \ac{DQN}-backed contextual bandit
& lab plus real classroom, n=271
& affect accuracy 82 to 88\%; about +15 points learning gain for mid-ability learners \\

\makecell[l]{Gong\\ (2025)\\ \cite{gong_design_2025}}
& multimodal: webcam, mic, text
& basic emotions (happy, sad, angry, neutral, surprised)
& \ac{DQN} (\ac{MDP})
& public datasets only; no real students
& fusion accuracy 0.91, AUC 0.94 (dataset-only, basic emotions) \\

\end{longtable}
\end{landscape}
\restoregeometry
\clearpage

\section{Research Gap}\label{sec:research_gap}

Facial expression recognition for education has matured from handcrafted features and shallow classifiers into deep temporal and transformer-based models that run on an ordinary webcam in real time \cite{lek_academic_2023,su_leveraging_2024,aly_comprehensive_2025,liu_dynamic_2026}. Public benchmarks such as \ac{DAiSEE} made the academic emotions of boredom, confusion, engagement and frustration a shared target \cite{gupta_daisee_2016}. Two limits still hold. Recognition accuracy on the four-level engagement task stays in the low-to-mid sixties even for strong models, and visually similar states such as confusion and frustration are still confused \cite{abedi_improving_2021,rianto_engagement_2025,lek_academic_2023}. More important for this thesis, almost all of this work stops at classification, an engagement index, or a dashboard, and does not feed its output into a tutoring decision \cite{rao_recognition_2021,komaravalli_detecting_2023,leong_deep_2020,aly_comprehensive_2025}. Facial expression recognition has therefore solved a perception problem without closing the loop to action.

In parallel, reinforcement learning has given tutoring systems a principled way to choose what to do next. Adaptation has been framed as a sequential decision problem, and policies have been learned for exercise selection, review scheduling, hints, presentation and learning-path sequencing \cite{azoulay_adaptive_2021,kubotani_rltutor_2021,sayed_ai-based_2023,fu_integrating_2025,zou_deep_2026}. Because online exploration on real learners is costly and slow, and because reward functions are easy to misspecify, this work has settled on simulation with synthetic students as the standard initial training stage \cite{dulacarnold_challenges_2021,sanz_ausin_leveraging_2019,li_reinforcement_2022,olukola_pedagogical_2026}. These tutors are strong on the cognitive side, but they represent the learner almost only through correctness, response time, hint use and knowledge-tracing estimates \cite{azoulay_adaptive_2021,fu_integrating_2025,guzman_developing_2025,zou_deep_2026}. The affective state of the learner does not enter the decision policy.

Emotion-aware tutoring has tried to bring affect back into the loop. The idea is more than fifteen years old, from multi-sensor classroom rigs \cite{woolf_affect-aware_2009} to recent systems that join emotion recognition with reinforcement learning \cite{perez_emotions_2024,govea_implementation_2024,kong_real-time_2025,gong_design_2025}. A field-level scoping review still reports that these systems recognise affect more reliably than they act on it, with scaffolding the most common and often hand-coded response \cite{fernandez-herrero_evaluating_2024}. Three problems recur. The strongest results depend on multimodal or physiological sensing, such as cameras combined with voice, posture, skin conductance or other biosignals, which is hard to reproduce in ordinary classrooms or at scale \cite{woolf_affect-aware_2009,standen_evaluation_2020,govea_implementation_2024,kong_real-time_2025}. Several systems detect only the basic emotions rather than the academic emotions that matter for learning \cite{govea_implementation_2024,gong_design_2025}. And the affect-to-action mapping is often fixed by hand or organised around a single control such as difficulty or tone \cite{megahed_modeling_2020,perez_emotions_2024,gutierrez_development_2025,sayed_ai-based_2023}.

Read together, these three lines leave a specific problem open. Affect can be sensed from an ordinary webcam, and tutoring actions can be learned in simulation, but the two have not been joined in a way that is at once deployable, learnable and testable. The affective signal is rarely part of the learned state and reward of the policy, since it is added as a rule or used only for detection \cite{fernandez-herrero_evaluating_2024,megahed_modeling_2020}. The affective controllers that do perform well rely on sensing that does not scale, which leaves the trade-off between sensing richness and deployment feasibility unresolved \cite{petrovica_emotion_2016,kong_real-time_2025,govea_implementation_2024}. The contribution of affect is also seldom isolated: the closest system to this thesis combines facial valence and arousal with Q-learning, but it draws affect from a static corpus rather than a live webcam \cite{perez_emotions_2024}. Reproducibility makes the problem harder, because the state, reward, transition assumptions and hyperparameters of these controllers are often left unspecified \cite{govea_implementation_2024,kong_real-time_2025}.

This open problem matters for more than engineering. The evidence that emotion-aware adaptation improves learning is still weak. The most carefully evaluated multi-sensor study found more engagement and less boredom under affect-aware adaptation, but no significant difference in learning gains \cite{standen_evaluation_2020}. At the same time, placing affect inside the reward is risky, because a policy that optimises an engagement proxy can reach a high reward while producing little mastery \cite{olukola_pedagogical_2026}. Deciding whether a webcam-derived affective signal genuinely changes a tutoring policy, and whether that change helps learning or only inflates the reward, therefore needs a controlled comparison under matched conditions. Without an explicit and reproducible protocol that isolates the affective channel, the question of whether sensed emotion should drive instruction cannot be answered.

The present thesis addresses this problem directly. It uses facial expression recognition from a single webcam, trained for the academic emotions of boredom, confusion, engagement and frustration, following the sensor-lite path argued for ordinary classrooms \cite{petrovica_emotion_2016,lek_academic_2023}. It places the resulting affective state inside the state of a reinforcement-learning tutor that selects from a discrete set of pedagogical actions, rather than a single difficulty rung \cite{azoulay_adaptive_2021,fu_integrating_2025}. It then evaluates the policy in a synthetic student environment with a stated state, reward and transition model, comparing an emotion-aware policy against a matched emotion-free baseline under the same student population and hyperparameters \cite{sanz_ausin_leveraging_2019,olukola_pedagogical_2026,weitekamp_tutorgym_2025}. This design keeps sensing within reach of real classrooms while integrating affective information directly into the reinforcement-learning state representation rather than relying on fixed pedagogical rules. Building on the research gaps identified in this chapter, the proposed framework combines a facial expression recognition module with an adaptive tutoring system to create a complete emotion-aware learning environment. The thesis also proposes a lightweight facial expression recognition architecture for academic emotion detection that achieves strong classification performance and serves as the affective input to the adaptive tutoring system. The next chapter describes the \ac{FER} architecture, the \ac{DQN} reinforcement-learning tutor, the full Markov decision process, the synthetic student environment, and the experimental framework used to evaluate the proposed emotion-aware tutoring framework under controlled simulation conditions.

